% 第5章 理想不可压缩流体的二维无旋和有旋流动 —— 依据原书扫描页手工精修版
\chapter{理想不可压缩流体的二维无旋和有旋流动}

本章应用理想流体动力学理论求解不可压缩流体的平面和轴对称流动，包括无旋流动和有旋流动.

\section*{5.1  不可压缩平面流动和轴对称流动的流函数}\addcontentsline{toc}{section}{5.1 不可压缩平面流动和轴对称流动的流函数}

\subsection*{1. 平面流动和轴对称流动的定义}

\noindent\textbf{定义5.1}\quad 流体质点在平行平面上运动，并且每一平面上流动都相同的流场称为平面平行流动，简称平面流动.

如果流动平行于$xy$平面，则平面流动速度场可表示为
\begin{equation*}
U_x=u(x,y,t)\,,\qquad U_y=v(x,y,t)\,,\qquad U_z=0
\tag{5.1}
\end{equation*}
理论上均匀来流垂直于无限长柱体的绕流是二维流动.实际流动中有不少情况可以近似为二维流动.例如，均匀来流垂直于长柱体的绕流的情况，见图5.1，假定柱体两端对柱体中部绕流的影响很小，因此柱体中部的绕流可近似为二维绕流.水平飞行的机翼绕流，机翼中部的流动，也可以近似为平面平行的二维流动.对于平行平面流动，只需要研究一个流动平面上的流动状态，如果不加说明，垂直于流动方向的长度取作单位长度.

\noindent\textbf{定义5.2}\quad 流体质点在通过固定轴线的子午面上运动，并且所有子午面上运动都相同的流场，称为轴对称流.

设流动的对称轴是$x$轴，则在圆柱坐标中轴对称流的速度场可表示为
\begin{equation*}
U_r=v_r(r,x,t)\,,\qquad U_x=v_x(r,x,t)\,,\qquad U_\theta=0
\tag{5.2}
\end{equation*}
在实际情况中有不少轴对称流动.回转体沿轴向运动驱动的流场或均匀来流平行于回转体轴线的绕流都属于轴对称流，例如：炮弹、火箭等无攻角飞行，见图5.2.

轴对称流动也可用球坐标来描述：
\begin{equation*}
U_R=v_R(R,\theta,t)\,,\qquad U_\theta=v_\theta(R,\theta,t)\,,\qquad U_\varphi=0
\tag{5.3}
\end{equation*}
二维流动只有两个空间变量，无论在分析或在数值计算方面都较一般三维流动要简单得多.凡是可能用二维近似的流动，首先做二维流动分析，这是一种比较简捷的方法.下面分别讲述轴对称和平面平行流动的分析方法.

\subsection*{2. 平面和轴对称流动的流函数}

不可压缩流体二维运动常用流函数$\Psi(x,y)$来表示流场和计算流动速度，流函数由连续方程导出.

\textbf{(1)平面流动的流函数}

不可压缩平面流动的连续方程是
\[
\frac{\partial u}{\partial x}+\frac{\partial v}{\partial y}=0
\]
利用斯托克斯公式，满足上述连续方程的$u,v$可以构成全微分：
\[
\mathrm{d}\Psi=u\,\mathrm{d}y-v\,\mathrm{d}x
\]
于是平面不可压缩流中的流函数可如下定义：

\noindent\textbf{定义5.3}\quad $\Psi=\displaystyle\int( u\,\mathrm{d}y-v\,\mathrm{d}x)$称为不可压缩平面流动的流函数.

对流函数求导数，得到它和速度之间的关系
\begin{equation*}
u=\frac{\partial\Psi}{\partial y}\,,\qquad v=-\frac{\partial\Psi}{\partial x}
\tag{5.4}
\end{equation*}
很容易证明式(5.4)满足不可压缩平面流动的连续方程.

\textbf{(2)轴对称流动的流函数}

柱坐标中不可压缩轴对称流动的连续方程可写成
\begin{equation*}
\frac{\partial(rv_r)}{\partial r}+\frac{\partial(rv_x)}{\partial x}=0
\tag{5.5}
\end{equation*}
根据斯托克斯定理，满足式(5.5)的两函数$rv_r,rv_x$可以构成全微分：
\[
\mathrm{d}\Psi=rv_r\,\mathrm{d}r-rv_x\,\mathrm{d}x
\]
于是有定义：

\noindent\textbf{定义5.4}\quad $\Psi=\displaystyle\int( rv_r\,\mathrm{d}r-rv_x\,\mathrm{d}x)$称为不可压缩轴对称流动的流函数.

对流函数求导数，可得它和速度场之间的关系：
\begin{equation*}
v_r=\frac{1}{r}\,\frac{\partial\Psi}{\partial r}\,,\qquad v_x=-\frac{1}{r}\,\frac{\partial\Psi}{\partial x}
\tag{5.6}
\end{equation*}
在数学运算上，流函数是将满足连续方程的两个函数$(u,v)$或$(v_r,v_x)$简化为一个函数$\Psi$.流函数的存在是以连续方程为前提，因此，用流函数描述流场，连续方程自然满足.

\textbf{(3)流函数的意义和性质}

\textcircled{1} \textbf{流函数的等值线是流线.}

以轴对称流动为例说明它的意义和证明它的性质，读者可以用相同的方法证明平面流动的流函数具有类似的性质.

\textbf{证明}\quad 由流线定义：
\[
\frac{\mathrm{d}x}{v_x}=\frac{\mathrm{d}r}{v_r}\qquad\text{或}\qquad v_x\,\mathrm{d}r-v_r\,\mathrm{d}x=0
\]
用流函数替代$v_x,v_r$得以下的流线方程：
\[
r\,\frac{\partial\Psi}{\partial r}\,\mathrm{d}r+r\,\frac{\partial\Psi}{\partial x}\,\mathrm{d}x=0
\]
消去$r$后上式成为全微分方程：
\[
\mathrm{d}\Psi=0
\]
于是证明了沿流线流函数为常数，即
\[
\Psi(x,r)=\mathrm{const.}
\]
在柱坐标中它是子午面上的流线，也是轴对称流中以流线为母线的回转流面方程.

用相同的方法可以证明
\[
\Psi(x,y)=\mathrm{const.}
\]
是不可压缩平面流动的流线，也是垂直于流动平面的流面方程（读者自己证明）.

\textcircled{2} \textbf{子午面上两流线间流函数值之差等于通过相应旋转流面间的体积流量除以$2\pi$.}

\textbf{证明}\quad 在$\Psi_1,\Psi_2$两流线之间任意联线$AB$，并以此联线作为母线作回转面，由于轴对称性，通过该回转面的体积流量为
\[
Q=\int_A^B 2\pi r\,\boldsymbol{U}\cdot\boldsymbol{n}\,\mathrm{d}s=2\pi\int_A^B r\,(v_x n_x+v_r n_r)\,\mathrm{d}s
\]
式中$\mathrm{d}s$是$AB$联线的微元弧长，$\boldsymbol{n}$为微元弧长的法向量，设联线的几何方程为：$r=r(s)$，$x=x(s)$，则微元弧线的切向量为：$s_x=\partial x/\partial s$，$s_r=\partial r/\partial s$，曲线的法向量和切向量之间有以下关系式：$n_x=s_r=\partial r/\partial s$，$n_r=-s_x=-\partial x/\partial s$.将它们和流函数表达式代入流量积分公式，得
\[
Q=2\pi\int_A^B r\,\Big(\frac{\partial\Psi}{r\,\partial r}\,\mathrm{d}r+\frac{\partial\Psi}{r\,\partial x}\,\mathrm{d}x\Big)=2\pi\int_A^B \mathrm{d}\Psi=2\pi\,(\Psi_B-\Psi_A)
\]
或
\[
\Psi_B-\Psi_A=\frac{Q}{2\pi}
\]
证毕.

在不可压缩的平面流动中，$\Psi_B-\Psi_A=q$，$q$是通过$A$和$B$间单位厚度（垂直于流动平面方向）的流量.仿效以上的证明，读者可以自己证实该性质.

\textcircled{3} \textbf{流函数等值线和势函数等值线正交.}

\textbf{证明}\quad 如果轴对称流动是无旋的，则流动还存在速度势$\Phi$，流场速度用速度势表示为
\[
v_x=\frac{\partial\Phi}{\partial x}\,,\qquad v_r=\frac{\partial\Phi}{\partial r}
\]
与轴对称速度场的流函数表达式(5.6)对比，可以导出轴对称无旋流动中流函数和势函数的关系式如下：
\begin{equation*}
v_x=\frac{\partial\Phi}{\partial x}=\frac{1}{r}\,\frac{\partial\Psi}{\partial r}\,,\qquad v_r=\frac{\partial\Phi}{\partial r}=-\frac{1}{r}\,\frac{\partial\Psi}{\partial x}
\tag{5.7}
\end{equation*}
利用式(5.7)就很容易证明流函数和势函数等值线的正交性.流函数和势函数等值线的法向量分别为：$\boldsymbol{n}_\psi=\dfrac{\nabla\psi}{|\nabla\psi|}$，$\boldsymbol{n}_\phi=\dfrac{\nabla\Phi}{|\nabla\Phi|}$，如$\boldsymbol{n}_\phi\cdot\boldsymbol{n}_\psi=0$，则两族等值线正交.由流函数和势函数的定义：
\[
\nabla\Phi=\frac{\partial\Phi}{\partial r}\,\boldsymbol{e}_r+\frac{\partial\Phi}{\partial x}\,\boldsymbol{e}_x=v_r\boldsymbol{e}_r+v_x\boldsymbol{e}_x
\]
\[
\nabla\Psi=\frac{\partial\Psi}{\partial r}\,\boldsymbol{e}_r+\frac{\partial\Psi}{\partial x}\,\boldsymbol{e}_x=rv_x\boldsymbol{e}_r-rv_r\boldsymbol{e}_x
\]
于是有
\[
\boldsymbol{n}_\phi\cdot\boldsymbol{n}_\psi=\frac{\nabla\Phi\cdot\nabla\Psi}{|\nabla\Phi|\,|\nabla\Psi|}=rv_r v_x-rv_x v_r=0
\]
这就证明了流函数等值线和势函数等值线正交.从几何上说，$\Psi$的等值线是流线，即$\Psi$等值线的切线方向是当地的速度方向；另一方面，无旋流动的速度$\boldsymbol{U}=\nabla\Phi$，它在速度势等值面的法线方向，即流体质点速度必垂直于等势面，于是从几何上也可证明流函数$\Psi$的等值线和势函数$\Phi$的等值线正交.

同理可以证明，不可压缩平面流动中也有同样性质.

\section*{5.2  不可压缩轴对称定常无旋流动}\addcontentsline{toc}{section}{5.2 不可压缩轴对称定常无旋流动}

\subsection*{1. 不可压缩轴对称无旋流动的流函数方程和定解条件}

\textbf{(1)不可压缩轴对称无旋流动的流函数方程}

流函数$\Psi$满足不可压缩连续方程，令它满足无旋条件，就得到不可压缩无旋流动的流函数方程.用柱坐标表示轴对称流中的旋度为
\[
\omega_\theta=\frac{\partial v_r}{\partial x}-\frac{\partial v_x}{\partial r}\,,\qquad \omega_r=\omega_z=0
\]
即轴对称流动中只有一个垂直于子午面的旋度分量，将流函数$v_x=\dfrac{1}{r}\dfrac{\partial\Psi}{\partial r}$，$v_r=-\dfrac{1}{r}\dfrac{\partial\Psi}{\partial x}$代入后得
\[
\frac{\partial}{\partial r}\Big(\frac{1}{r}\,\frac{\partial\Psi}{\partial r}\Big)+\frac{\partial}{\partial x}\Big(\frac{1}{r}\,\frac{\partial\Psi}{\partial x}\Big)=0
\]
上式可展开为
\begin{equation*}
\frac{\partial^{2}\Psi}{\partial x^{2}}+\frac{\partial^{2}\Psi}{\partial r^{2}}-\frac{1}{r}\,\frac{\partial\Psi}{\partial r}=0
\tag{5.8}
\end{equation*}
式(5.8)表明：不可压缩轴对称无旋流动的流函数方程也是二阶线性椭圆型偏微分方程，但它不是拉普拉斯方程.

\textbf{(2)用流函数表示的绕流问题边界条件}

\textcircled{1} 物面条件：定常绕流问题中轴对称物面是回转流面，而流函数在轴对称流面上是常数，因此用流函数表示的物面不可穿透条件为

物面$\Sigma$上：
\begin{equation*}
\Psi=\mathrm{const.}
\tag{5.9}
\end{equation*}

\textcircled{2} 无穷远处来流条件是（来流平行于$x$轴）：
\[
\boldsymbol{U}=U_\infty\,\boldsymbol{e}_x
\]
用流函数表示时，应是
\begin{equation*}
\lim_{x\to\infty}\frac{1}{r}\,\frac{\partial\Psi}{\partial r}=U_\infty
\tag{5.10a}
\end{equation*}
上式对$r$积分，可表示为
\begin{equation*}
\lim_{x\to\infty}\Psi=\frac{1}{2}\,U_\infty r^{2}
\tag{5.10b}
\end{equation*}

\textcircled{3} 轴对称流函数在球坐标系中的表达式

轴对称流也可用球坐标表示.
\[
U_R=v_R(R,\theta,t)\,,\qquad U_\theta=v_\theta(R,\theta,t)
\]
球坐标中轴对称流动的连续性方程为
\begin{equation*}
\frac{\partial(R^{2}\sin\theta\,v_R)}{\partial R}+\frac{\partial(R\sin\theta\,v_\theta)}{\partial\theta}=0
\tag{5.11}
\end{equation*}
仍然应用斯托克斯定理，满足式(5.11)的$R^{2}\sin\theta v_R$和$R\sin\theta v_\theta$可用一流函数$\Psi$替代，它在球坐标系中的定义为
\[
\frac{\partial\Psi}{\partial\theta}=R^{2}\sin\theta\,v_R\,,\qquad \frac{\partial\Psi}{\partial R}=-R\sin\theta\,v_\theta
\]
或
\begin{equation*}
v_R=\frac{1}{R^{2}\sin\theta}\,\frac{\partial\Psi}{\partial\theta}\,,\qquad v_\theta=\frac{-1}{R\sin\theta}\,\frac{\partial\Psi}{\partial R}
\tag{5.12}
\end{equation*}
在球坐标中，轴对称流动的旋度表达式是
\[
\nabla\times\boldsymbol{U}=\Big(\frac{\partial v_\theta}{\partial R}+\frac{v_\theta}{R}-\frac{1}{R}\,\frac{\partial v_R}{\partial\theta}\Big)\boldsymbol{e}_\varphi
\]
令旋度等于零，并代入速度的流函数表达式，很容易导出球坐标中不可压缩轴对称无旋流的流函数方程为
\begin{equation*}
\frac{\partial^{2}\Psi}{\partial R^{2}}+\frac{1}{R^{2}}\,\frac{\partial^{2}\Psi}{\partial\theta^{2}}-\frac{\cot\theta}{R^{2}}\,\frac{\partial\Psi}{\partial\theta}=0
\tag{5.13}
\end{equation*}
引入流函数后，不可压缩轴对称无旋流动也可以由流函数方程及其定解条件解出.读者不禁要问：既然有速度势方程及其定解条件可以来解不可压缩无旋流动，为何还要用流函数方程来解呢？理由是：①用流函数求解流动时，物面上所用的是第一类边界条件，而用速度势求解无旋流动时物面上是第二类边界条件，当用基本解叠加或数值方法求解时，第一类边界条件问题较为容易求解.②求得流函数后，直接可得到流线谱，这就可以比较直观地来描述流动，从而也可以比较方便地判断解的合理性和正确性.③对于不可压缩流体的二维流动，不论是定常的还是非定常的，也不论流动是理想无旋的还是真实有粘性的有旋流动，流函数总是存在的（它是连续方程的一种简化式），因此有粘性的不可压缩流体二维流动也常用流函数方法来分析和求解.此外，可压缩流体的二维定常流动也存在流函数，可以用流函数方法求解（请读者完成本章习题5.18的证明）.

\subsection*{2. 流函数表示的基本解和叠加法}

第4章给出的各种无旋流动基本解的速度势表达式，可通过流函数与势函数的变换关系$\dfrac{\partial\Phi}{\partial x}=\dfrac{1}{r}\dfrac{\partial\Psi}{\partial r}$，$\dfrac{\partial\Phi}{\partial r}=-\dfrac{1}{r}\dfrac{\partial\Psi}{\partial x}$（式(5.7)）求出轴对称基本解的流函数表达式：
\[
\Psi=\int\Big(\frac{\partial\Psi}{\partial r}\,\mathrm{d}r+\frac{\partial\Psi}{\partial x}\,\mathrm{d}x\Big)=\int\Big(r\,\frac{\partial\Phi}{\partial x}\,\mathrm{d}r-r\,\frac{\partial\Phi}{\partial r}\,\mathrm{d}x\Big)
\]
已知任意一个不可压缩轴对称无旋流动的速度势$\Phi$，由上式通过积分可计算流函数$\Psi$.下面写出若干不可压缩无旋流动基本解流函数表达式，请读者自己用积分式推导证明.

\textbf{(1)均匀流的流函数表示式}
\begin{equation*}
\Psi=\frac{1}{2}\,U_\infty R^{2}\sin^{2}\theta\qquad(\text{球坐标})
\tag{5.14a}
\end{equation*}
\begin{equation*}
\Psi=\frac{1}{2}\,U_\infty r^{2}\qquad(\text{柱坐标})
\tag{5.14b}
\end{equation*}

\textbf{(2)点源的流函数表示式}
\[
\Psi=\frac{Q}{4\pi}\,(1-\cos\theta)\qquad(\text{球坐标})
\]
流函数的积分常数不影响速度场，因此点源的流函数也可简写为
\begin{equation*}
\Psi=-\frac{Q}{4\pi}\cos\theta\qquad(\text{球坐标})
\tag{5.15a}
\end{equation*}
柱坐标中点源的流函数为
\[
\Psi=\frac{Q}{4\pi}\Big(1-\frac{x}{\sqrt{x^{2}+r^{2}}}\Big)\qquad(\text{柱坐标})
\]
它也可简化为
\begin{equation*}
\Psi=-\frac{Q}{4\pi}\,\frac{x}{\sqrt{x^{2}+r^{2}}}\qquad(\text{柱坐标})
\tag{5.15b}
\end{equation*}

\textbf{(3)偶极子的流函数表达式}
\begin{equation*}
\Psi=\frac{m}{4\pi}\,\frac{\sin^{2}\theta}{R}\qquad(\text{球坐标})
\tag{5.16a}
\end{equation*}
\begin{equation*}
\Psi=\frac{m}{4\pi}\,\frac{r^{2}}{(x^{2}+r^{2})^{3/2}}\qquad(\text{柱坐标})
\tag{5.16b}
\end{equation*}
不可压缩流体轴对称无旋流动的流函数方程(5.8)或(5.13)是线性方程，因此也可以用基本解的线性叠加方法求这种流动的流函数解.下面利用流函数基本解求解一个轴对称绕流问题.

\textbf{(4)不可压缩流体绕兰金回转体的定常无旋流动}

\noindent\textbf{例5.1}\quad 设在$x=d$上有一点源，其强度为$Q$；在$x=-d$上有一点汇，其强度为$-Q$.将它们和均匀流$\boldsymbol{U}=-U_\infty\boldsymbol{i}$叠加，求叠加流场的流函数和绕流的物面型线.

\noindent\textbf{解}\quad 叠加的流函数为
\[
\Psi=-\frac{1}{2}\,U_\infty r^{2}+\frac{1}{4\pi}\,\frac{Q(x+d)}{\sqrt{(x+d)^{2}+r^{2}}}-\frac{1}{4\pi}\,\frac{Q(x-d)}{\sqrt{(x-d)^{2}+r^{2}}}
\]
令$\Psi(x,r)=0$，得流线方程为
\[
-\frac{1}{2}\,U_\infty r^{2}+\frac{1}{4\pi}\,\frac{Q(x+d)}{\sqrt{(x+d)^{2}+r^{2}}}-\frac{1}{4\pi}\,\frac{Q(x-d)}{\sqrt{(x-d)^{2}+r^{2}}}=0
\]
$r=0$满足流线方程，就是说$\Psi(x,r)=0$的流线通过$x$轴.另外，$\Psi(x,r)=0$的流线还有解
\[
r^{2}-\frac{b^{2}(x+d)}{\sqrt{(x+d)^{2}+r^{2}}}+\frac{b^{2}(x-d)}{\sqrt{(x-d)^{2}+r^{2}}}=0
\]
式中$b^{2}=Q/(2\pi U_\infty)$.因为$\Big|\dfrac{(x+d)}{\sqrt{(x+d)^{2}+r^{2}}}\Big|$和$\Big|\dfrac{(x-d)}{\sqrt{(x-d)^{2}+r^{2}}}\Big|$都小于1，因此$r^{2}<2b^{2}$，就是说$\Psi(x,r)=0$还有一封闭分支流线.轴对称流动的流线是空间回转曲面和坐标面的交线，因此上述封闭分支流线代表封闭的卵形回转面.于是，均匀绕流和一对等强度点源和点汇叠加，得到绕卵形回转体的不可压缩无旋绕流，该回转体又称兰金（Rankine）物体（见图5.4）.通过计算绕流物面的驻点（速度等于零的点）可以计算该物体的长度.由于兰金物体对$y$轴是对称的，可令驻点坐标为$(A,0)$，$(-A,0)$.流动速度可由流函数求导数获得，对$x$求导数，它等于$v_r$，并令它等于零，得到以下方程：
\[
(A^{2}-d^{2})=2b^{2}Ad
\]
当给定来流速度$U_\infty$和源强$Q$以及点源和点汇间的距离$2d$，由上式可以求出相应的兰金物体的长度$2A$.在物面方程中，令$x=0$，还可以求出兰金物体的宽度$H$满足的代数方程：
\[
H^{2}=2b^{2}d\big/\sqrt{H^{2}+d^{2}}
\]
令流函数等于任意常数，还可以求出流动的流线谱，如图5.4所示.求出流函数后，由流函数定义式(5.6)可以求出速度场，然后，由伯努利公式计算流场和物面上的压强系数.与圆球绕流情况相仿，理想不可压缩流体绕卵形物体的无旋流动的合力等于零.与实际情况相比，卵形体前部的理论压强系数和实际情况符合良好，物体后部附近流动发生分离，实际压强远远小于理论压强，因而产生卵形体上的流动阻力.

\noindent\textbf{例5.2}\quad 设回转体的物面方程为$r=f(x)$，并有$f(0)=f(L)=0$，$f'(0)<\infty$，$f'(L)<\infty$，即物面两端为尖缘、长为$L$的回转体.求不可压缩理想流体绕该回转体的无旋流动.

\noindent\textbf{解}\quad 仍用奇点叠加法来求解该问题，现在需要求出在轴线上分布何种奇点，才能生成绕给定物面的无旋流动.设分布点源$q(x)$，在长$\mathrm{d}\xi$线段上源强度为$q(\xi)\,\mathrm{d}\xi$，它的流函数为
\[
\mathrm{d}\Psi=\frac{q(\xi)\,\mathrm{d}\xi}{4\pi}\bigg[1-\frac{x-\xi}{\sqrt{(x-\xi)^{2}+r^{2}}}\bigg]
\]
分布点源与均匀来流（沿$x$轴正方向）叠加的流函数为
\[
\Psi(x,r)=\frac{1}{2}\,U_\infty r^{2}+\frac{1}{4\pi}\int_0^{L}q(\xi)\bigg[1-\frac{x-\xi}{\sqrt{(x-\xi)^{2}+r^{2}}}\bigg]\mathrm{d}\xi
\]
物面型线是封闭周线，因而流出物面的总体积流量应当等于零：
\[
\oint v_n\,\mathrm{d}S=0
\]
这要求在周线内点源的总流量为零，即$\displaystyle\int_0^{L}q(\xi)\,\mathrm{d}\xi=0$，故得
\begin{equation*}
\Psi(x,r)=\frac{1}{2}\,U_\infty r^{2}-\frac{1}{4\pi}\int_0^{L}\frac{q(\xi)\,(x-\xi)}{\sqrt{(x-\xi)^{2}+r^{2}}}\,\mathrm{d}\xi
\tag{5.17}
\end{equation*}
要求物面是流线，而且该流线与$r=0$在同一流线上，因而必须有
\[
(\Psi)_{r=0}=(\Psi)_{r=f(x)}
\]
将以上关系代入式(5.17)，得确定点源强度的积分方程：
\begin{equation*}
U_\infty f^{2}(x)-\frac{1}{2\pi}\int_0^{L}\frac{q(\xi)\,(x-\xi)}{\sqrt{(x-\xi)^{2}+f^{2}(x)}}\,\mathrm{d}\xi=0
\tag{5.18}
\end{equation*}
式中$f(x)$是物面已知函数.给定$f(x)$后，由积分方程反演可求出点源分布.对于任意物面函数，一般不可能求出解析解.常常用数值积分法离散式(5.18)，令$q_j=q(\xi_j)$，$f_i=f(x_i)$，式(5.18)可写成
\begin{equation*}
C_{ij}\,q_j=U_\infty f_i^{2}\,,\qquad i=1,N\,,\qquad j=1,N
\tag{5.19}
\end{equation*}
式中
\begin{equation*}
C_{ij}=\frac{1}{2\pi}\,\frac{(x_i-\xi_j)\,\Delta L}{\sqrt{(x_i-\xi_j)^{2}+f_i^{2}}}
\tag{5.20}
\end{equation*}
$\Delta L=L/N$为离散步长，用矩阵求逆便可得到$q_j$的分布.求出$q_j$分布后代入流函数公式，求出流函数，就可以得到速度场解.

要计算绕流流场的压强分布，可用伯努利公式.

\section*{5.3  解不可压缩平面无旋流动问题的复变函数方法}\addcontentsline{toc}{section}{5.3 解不可压缩平面无旋流动问题的复变函数方法}

\subsection*{1. 不可压缩平面无旋流动速度势问题的提法}

由于垂直于流动平面的速度$w=\partial\Phi/\partial z\equiv 0$，不可压缩平面无旋流动的速度势方程为
\begin{equation*}
\frac{\partial^{2}\Phi}{\partial x^{2}}+\frac{\partial^{2}\Phi}{\partial y^{2}}=0
\tag{5.21}
\end{equation*}
给定流动的速度势边界条件，求解二维拉普拉斯方程，就可得到流动的势函数解，其绕流问题的定解条件是

（1）来流条件是二维均匀流，应用上一章均匀来流的势函数表达式：
\begin{equation*}
|\boldsymbol{x}|\to\infty:\ \Phi=U_\infty x+V_\infty y
\tag{5.22}
\end{equation*}
（2）固壁物面周线$L$上的不可穿透条件$\boldsymbol{n}\cdot\nabla\Phi=0$可以写成
\begin{equation*}
L:F(x,y)=0\ \text{上}:\ \frac{\partial\Phi}{\partial x}\,\frac{\partial F}{\partial x}+\frac{\partial\Phi}{\partial y}\,\frac{\partial F}{\partial y}=0
\tag{5.23}
\end{equation*}
式中$F(x,y)=0$是物面型线方程.

（3）绕封闭物面周线的二维流场是双连通域，因此还需给出绕周线上的环量，前面已经证明，绕不可缩周线上的速度环量等于速度势的增量，因此环量条件可表示为
\begin{equation*}
\oint_L\boldsymbol{U}\cdot\mathrm{d}\boldsymbol{x}=\oint\mathrm{d}\Phi=\Gamma
\tag{5.24}
\end{equation*}

\subsection*{2. 不可压缩平面流动的流函数及平面无旋流动的流函数问题提法}

\textbf{(1)平面不可压缩无旋绕流问题的流函数方程}

平面流动的涡量场为
\begin{equation*}
\omega_z=\frac{\partial v}{\partial x}-\frac{\partial u}{\partial y}\,,\qquad \omega_x=\omega_y=0
\tag{5.25}
\end{equation*}
即平面流动中只有一个垂直于流动平面的涡量分量，平面无旋流动条件可写成
\begin{equation*}
\frac{\partial v}{\partial x}-\frac{\partial u}{\partial y}=0
\tag{5.26}
\end{equation*}
将流函数关系式(5.4)：$u=\partial\Psi/\partial y$，$v=-\partial\Psi/\partial x$代入无旋条件，得
\begin{equation*}
\frac{\partial^{2}\Psi}{\partial x^{2}}+\frac{\partial^{2}\Psi}{\partial y^{2}}=0
\tag{5.27}
\end{equation*}
式(5.27)表明不可压缩平面无旋流动的流函数和势函数一样满足拉普拉斯方程，这一点和不可压缩轴对称无旋流动不同，轴对称流的势函数满足拉普拉斯方程而流函数并不满足拉普拉斯方程.给定流函数边界条件，求解拉普拉斯方程，也可得到平面无旋流动的流函数解.

\textbf{(2)平面绕流的流函数方程定解条件为}

\textcircled{1} 二维均匀来流条件可写作
\begin{equation*}
|\boldsymbol{x}|\to\infty\,,\qquad \frac{\partial\Psi}{\partial y}=U_\infty\,,\qquad \frac{\partial\Psi}{\partial x}=-V_\infty
\tag{5.28a}
\end{equation*}
或用均匀流动的流函数表达式：
\begin{equation*}
\Psi\big|_{|\boldsymbol{x}|\to\infty}=U_\infty y-V_\infty x
\tag{5.28b}
\end{equation*}
\textcircled{2} 物面的不可穿透条件，即物面是流线，因此物面上流函数是常数：
\begin{equation*}
L:F(x,y)=0\ \text{上}\,,\ \Psi(x,y)=\mathrm{const.}
\tag{5.29}
\end{equation*}
\textcircled{3} 绕物面的环量条件
\begin{equation*}
\oint\boldsymbol{U}\cdot\mathrm{d}\boldsymbol{x}=\int u\,\mathrm{d}x+v\,\mathrm{d}y=\oint_l\Big(\frac{\partial\Psi}{\partial y}\,\mathrm{d}x-\frac{\partial\Psi}{\partial x}\,\mathrm{d}y\Big)=\Gamma
\tag{5.30}
\end{equation*}

\subsection*{3. 平面不可压缩无旋流动的复势和复速度}

由于平面不可压缩无旋流动的势函数和流函数都满足拉普拉斯方程，它们之间又存在以下关系：
\begin{equation*}
u=\frac{\partial\Phi}{\partial x}=\frac{\partial\Psi}{\partial y}\,,\qquad v=\frac{\partial\Phi}{\partial y}=-\frac{\partial\Psi}{\partial x}
\tag{5.31}
\end{equation*}
满足上述关系的一对调和函数，称为共轭调和函数.在数学上，共轭调和函数可以用复变函数表示.下面应用复变函数的理论求解平面不可压缩无旋流动.

\textbf{(1)复势}

\noindent\textbf{定义5.5}\quad 以速度势为实部、流函数为虚部组成的复函数称为复势，常用$W(z)$表示：
\begin{equation*}
W(z)=\Phi(x,y)+\mathrm{i}\Psi(x,y)
\tag{5.32}
\end{equation*}
式中$\mathrm{i}=\sqrt{-1}$，$z=x+\mathrm{i}y=r\exp(\mathrm{i}\theta)$，$x,y$是平面直角坐标；$r,\theta$是平面极坐标的向径和辐角.

复势表达式中速度势和流函数都是调和函数，并有共轭关系式(5.31)，共轭关系式又称柯西-黎曼条件，复变函数理论已经证明：满足柯西-黎曼条件的一对共轭调和函数组成的复函数是解析函数，因此复势是解析函数.解析函数有以下性质.

\textcircled{1} 复势的方向导数和求导方向无关，它只是复变量$z$的一元函数.

实数自变量$x,y$和复数$z$之间有以下关系：
\[
x=\frac{z+\bar{z}}{2}\,,\qquad y=\frac{z-\bar{z}}{2\mathrm{i}}
\]
复数加上杠表示它的复共轭.由上式可以得到：$\dfrac{\partial x}{\partial z}=\dfrac{1}{2}$，$\dfrac{\partial x}{\partial\bar{z}}=\dfrac{1}{2}$，$\dfrac{\partial y}{\partial z}=\dfrac{1}{2\mathrm{i}}=-\dfrac{\mathrm{i}}{2}$，$\dfrac{\partial y}{\partial\bar{z}}=-\dfrac{1}{2\mathrm{i}}=\dfrac{\mathrm{i}}{2}$.

利用实数自变量$x,y$和复数$z$之间的关系，函数$\Phi,\Psi$也可以用自变量$z$和$\bar{z}$表示为$\Phi(z,\bar{z})$，$\Psi(z,\bar{z})$.复势的导数有
\[
\frac{\partial W}{\partial z}=\frac{\partial\Phi}{\partial z}+\mathrm{i}\,\frac{\partial\Psi}{\partial z}\,,\qquad \frac{\partial W}{\partial\bar{z}}=\frac{\partial\Phi}{\partial\bar{z}}+\mathrm{i}\,\frac{\partial\Psi}{\partial\bar{z}}
\]
现在要证明$\dfrac{\partial W}{\partial\bar{z}}=0$和$\dfrac{\partial W}{\partial z}=\dfrac{\partial W}{\partial x}=\dfrac{\partial W}{\mathrm{i}\partial y}$.第一个等式表明复势只是自变量$z$的函数；第二个等式表明复势的导数和方向无关.先证明第一个等式.
\begin{align*}
\frac{\partial W}{\partial\bar{z}}&=\frac{\partial\Phi}{\partial\bar{z}}+\mathrm{i}\,\frac{\partial\Psi}{\partial\bar{z}}=\frac{\partial\Phi}{\partial x}\,\frac{\partial x}{\partial\bar{z}}+\frac{\partial\Phi}{\partial y}\,\frac{\partial y}{\partial\bar{z}}+\mathrm{i}\Big(\frac{\partial\Psi}{\partial x}\,\frac{\partial x}{\partial\bar{z}}+\frac{\partial\Psi}{\partial y}\,\frac{\partial y}{\partial\bar{z}}\Big)\\
&=\frac{1}{2}\,\frac{\partial\Phi}{\partial x}+\frac{\mathrm{i}}{2}\,\frac{\partial\Phi}{\partial y}+\mathrm{i}\Big(\frac{1}{2}\,\frac{\partial\Psi}{\partial x}+\frac{\mathrm{i}}{2}\,\frac{\partial\Psi}{\partial y}\Big)=\frac{1}{2}\Big(\frac{\partial\Phi}{\partial x}-\frac{\partial\Psi}{\partial y}\Big)+\frac{\mathrm{i}}{2}\Big(\frac{\partial\Phi}{\partial y}+\frac{\partial\Psi}{\partial x}\Big)
\end{align*}
由于$\dfrac{\partial\Phi}{\partial x}=\dfrac{\partial\Psi}{\partial y}$和$\dfrac{\partial\Phi}{\partial y}=-\dfrac{\partial\Psi}{\partial x}$，于是证明了：
\[
\frac{\partial W}{\partial\bar{z}}=0
\]
进一步计算$\dfrac{\partial W}{\partial z}$.
\begin{align*}
\frac{\partial W}{\partial z}&=\frac{\partial\Phi}{\partial z}+\mathrm{i}\,\frac{\partial\Psi}{\partial z}=\frac{\partial\Phi}{\partial x}\,\frac{\partial x}{\partial z}+\frac{\partial\Phi}{\partial y}\,\frac{\partial y}{\partial z}+\mathrm{i}\Big(\frac{\partial\Psi}{\partial x}\,\frac{\partial x}{\partial z}+\frac{\partial\Psi}{\partial y}\,\frac{\partial y}{\partial z}\Big)\\
&=\frac{1}{2}\,\frac{\partial\Phi}{\partial x}+\frac{1}{2\mathrm{i}}\,\frac{\partial\Phi}{\partial y}+\mathrm{i}\Big(\frac{1}{2}\,\frac{\partial\Psi}{\partial x}+\frac{1}{2\mathrm{i}}\,\frac{\partial\Psi}{\partial y}\Big)=\frac{1}{2}\Big(\frac{\partial\Phi}{\partial x}+\frac{\partial\Psi}{\partial y}\Big)+\frac{\mathrm{i}}{2}\Big(-\frac{\partial\Phi}{\partial y}+\frac{\partial\Psi}{\partial x}\Big)
\end{align*}
由于$\dfrac{\partial\Phi}{\partial x}=\dfrac{\partial\Psi}{\partial y}$和$\dfrac{\partial\Phi}{\partial y}=-\dfrac{\partial\Psi}{\partial x}$，于是证明了
\[
\frac{\partial W}{\partial z}=\frac{\partial\Phi}{\partial x}+\mathrm{i}\,\frac{\partial\Psi}{\partial x}=\frac{\partial W}{\partial x}
\]
上式也可写成$\dfrac{\partial W}{\partial z}=\dfrac{\partial\Phi}{\mathrm{i}\partial y}+\dfrac{\partial\Psi}{\partial y}=\dfrac{\partial W}{\mathrm{i}\partial y}$，于是证明了复势的方向导数和求导方向无关.由于$\dfrac{\partial W}{\partial\bar{z}}=0$，可以把导数写成一元复函数的导数形式：
\[
W'(z)=\frac{\mathrm{d}W}{\mathrm{d}z}
\]
以上证明的复势性质说明复势是解析函数，它的意义在于把两个实自变量的二元函数$\Phi(x,y)+\mathrm{i}\Psi(x,y)$（由于调和函数的共轭性）组成复自变量$z$的一元函数.解析函数还有以下性质.

\textcircled{2} 解析函数的和是解析函数.

\textcircled{3} 解析函数的复合函数也是解析函数，即如$f(z)$是解析函数，$z=Z(\zeta)$也是解析函数，则$F(\zeta)=f[Z(\zeta)]$也是解析函数.

\textcircled{4} 解析函数的导数是解析函数，也就是说解析函数可以无限求导.

\textcircled{5} 解析函数的积分是解析函数.

\textcircled{6} 在全平面上处处解析的函数是常数.

\textcircled{7} 在不包含$z_0=0$点的有限域中，解析函数的一般展开式为
\begin{equation*}
f(z)=\sum_{n=-\infty}^{+\infty}a_n(z-z_0)^{n}
\tag{5.33}
\end{equation*}
称式(5.33)为罗朗（Laurent）级数.

\textcircled{8} 留数公式
\begin{equation*}
a_{-1}=\frac{1}{(m-1)!}\lim_{z\to z_0}\frac{\mathrm{d}^{m-1}}{\mathrm{d}z^{m-1}}\big[(z-z_0)^{m}f(z)\big]
\tag{5.34a}
\end{equation*}
\begin{equation*}
\oint_L f(z)\,\mathrm{d}z=2\pi\mathrm{i}\,a_{-1}
\tag{5.34b}
\end{equation*}
$L$为包含$z_0$的封闭周线，$z_0$是$f(z)$的$m$阶极点.

复势是解析函数，所以它也具有以上性质.

\textbf{(2)复速度}

\noindent\textbf{定义5.6}\quad 以平面无旋流场的速度分量组成的复数$U=u+\mathrm{i}v$称为复速度.

复速度与复势的关系由复势求导给出：
\[
\frac{\mathrm{d}W}{\mathrm{d}z}=\frac{\partial W}{\partial x}=\frac{\partial\Phi}{\partial x}+\mathrm{i}\,\frac{\partial\Psi}{\partial x}=u-\mathrm{i}v
\]
即复势的导数等于复速度的共轭：
\begin{equation*}
\frac{\mathrm{d}W}{\mathrm{d}z}=\bar{U}
\tag{5.35}
\end{equation*}
以上给出了平面不可压缩无旋流场的复变函数描述法，以及它们的主要性质，下面讨论应用复函数求解不可压缩流体无旋绕流问题的方法.

\subsection*{4. 不可压缩平面无旋绕流问题的复势提法}

由于解析函数的实部和虚部都是调和函数，因此给定某一解析函数，取它的实部和虚部都满足拉普拉斯方程，也就是说，给定一个解析函数就有一个与之相应的平面无旋流动.但是对于具体的绕流问题，速度势和流函数都必须满足特定的边界条件.这就是说，对具体的绕流问题，必须求出满足特定边界条件的复势.由速度势和流函数的边界条件可导出复势的边界条件如下.

\textbf{(1)物面型线$L$上的不可穿透条件}

前面已经导出用流函数表示的物面不可穿透条件是在物面上流函数应等于常数，也就是复势虚部等于常数，因而有
\begin{equation*}
L:F(x,y)=0\,,\qquad \Psi=\mathrm{Im}W(z)=\mathrm{const.}
\tag{5.36}
\end{equation*}
\textbf{(2)无穷远处来流条件}

无穷远处来流速度给定，也就是给定无穷远处复速度$U=U_\infty+\mathrm{i}V_\infty$，因而有
\begin{equation*}
\Big(\frac{\mathrm{d}W}{\mathrm{d}z}\Big)_{z\to\infty}=U_\infty-\mathrm{i}V_\infty
\tag{5.37}
\end{equation*}
\textbf{(3)不可缩周线上的环量}

有封闭周线的平面无旋绕流问题是双连通域中的有势流动，因此必须给定绕物面周线（即不可缩周线）的环量，它可用速度势的增量表示：
\[
\int_l\boldsymbol{U}\cdot\mathrm{d}\boldsymbol{x}=\oint_l\mathrm{d}\Phi=\Gamma
\]
用复势表示的环量条件导出如下：因为绕流物面上流函数是常量，见式(5.36)，因而$\displaystyle\oint_l\mathrm{d}\Psi=0$，从而绕封闭物面周线的复势增量等于绕物面的速度环量：
\begin{equation*}
\oint_l\mathrm{d}W=\oint_l\mathrm{d}\Phi=\Gamma
\tag{5.38}
\end{equation*}
以上建立了平面无旋绕流问题的复势边界条件.在复变函数理论中已经证明，满足边界条件式(5.36)、式(5.37)和式(5.38)的解析函数是唯一的.也就是说，满足边界条件的复势是唯一的，它对应于唯一的不可压缩无旋流动的速度场.

现在有了三种求解不可压缩平面无旋流动的方法.第一，速度势方程的边值问题；第二，流函数方程的边值问题；第三，复势的边值问题.前两种方法是求解实变量的二维拉普拉斯方程的边界值问题；第三种方法是已知边值寻求解析复变函数.前两种方法的原理在前面已介绍过，它们仍然可以应用到平面无旋流中.本节主要介绍用复变函数方法求解绕流问题，具体来说用奇点叠加方法和保角映像方法来求得满足边界条件的复势.

\subsection*{5. 平面不可压缩无旋流动的基本解及其复势表达式}

因为解析函数的和仍是解析函数，所以也可用奇点叠加的方法求解不可压缩平面无旋绕流问题的复势，下面先把不可压缩平面无旋流动的基本解用复势表示.

\textbf{(1)均匀流的复势}

均匀流的速度势为
\[
\Phi=U_\infty x+V_\infty y
\]
应用不可压缩无旋流动的流函数和势函数之间的共轭关系式(5.26)，可得流函数
\[
\Psi=U_\infty y-V_\infty x
\]
因而，均匀流的复势为
\begin{equation*}
W(z)=\Phi+\mathrm{i}\Psi=U_\infty z-\mathrm{i}V_\infty z=(U_\infty-\mathrm{i}V_\infty)z
\tag{5.39}
\end{equation*}
就是说，均匀流的复势是线性解析函数.

\textbf{(2)平面点源}

垂直于$xy$平面的无限长直线上有均匀分布的点源，其单位长度的点源强度为$q$，称这种分布为线源.线源以外的速度场可用第2章公式(2.52b)导出：
\[
\boldsymbol{U}_s=\frac{1}{4\pi}\iiint_D\theta(\xi,\eta,\zeta)\,\frac{\boldsymbol{R}}{R^{3}}\,\mathrm{d}\xi\,\mathrm{d}\eta\,\mathrm{d}\zeta
\]
设线源置于$z$轴，则：$\theta(\xi,\eta,\zeta)\,\mathrm{d}\xi\,\mathrm{d}\eta\,\mathrm{d}\zeta=q\,\mathrm{d}\zeta$，$\boldsymbol{R}=x\boldsymbol{i}+y\boldsymbol{j}+(z-\zeta)\boldsymbol{k}$，由线源产生的速度场
\[
\boldsymbol{U}=\frac{1}{4\pi}\int_{-\infty}^{+\infty}q\,[x\boldsymbol{i}+y\boldsymbol{j}+(z-\zeta)\boldsymbol{k}]\,\mathrm{d}\zeta\big/\sqrt{\big[x^{2}+y^{2}+(z-\zeta)^{2}\big]^{3}}
\]
已知积分等式：$\displaystyle\int_{-\infty}^{+\infty}\mathrm{d}\zeta\big/\sqrt{(a^{2}+\zeta^{2})^{3}}=2/a^{2}$和$\displaystyle\int_{-\infty}^{+\infty}\zeta\,\mathrm{d}\zeta\big/\sqrt{(a^{2}+\zeta^{2})^{3}}=0$，于是，线源的速度场为：$\boldsymbol{U}=\dfrac{q}{2\pi}\Big(\dfrac{x\boldsymbol{i}}{x^{2}+y^{2}}+\dfrac{y\boldsymbol{j}}{x^{2}+y^{2}}\Big)$.线源的速度场是有势的，它的势
\begin{equation*}
\Phi(x,y)=\int( u\,\mathrm{d}x+v\,\mathrm{d}y)=\frac{q}{2\pi}\ln\sqrt{x^{2}+y^{2}}=\frac{q}{2\pi}\ln r
\tag{5.40}
\end{equation*}
式中$r$是平面极坐标向径.由式(5.40)可见，无限长直线上均匀分布点源生成平面有势流场，这种平面流动简称为平面点源流动.由流函数和速度势关系式(5.31)，很容易算出不可压缩平面点源的流函数如下：
\begin{equation*}
\Psi=\frac{q}{2\pi}\arctan\frac{y}{x}=\frac{q}{2\pi}\,\theta
\tag{5.41}
\end{equation*}
式(5.41)中$\theta$是平面极坐标的辐角.代入复势公式，点源的复势表达式为
\begin{equation*}
W(z)=\Phi+\mathrm{i}\Psi=\frac{q}{2\pi}\,(\ln r+\mathrm{i}\theta)=\frac{q}{2\pi}\ln z
\tag{5.42a}
\end{equation*}
位于任意点$z_0$的平面点源复势可以用坐标平移的方法求出为
\begin{equation*}
W(z)=\Phi+\mathrm{i}\Psi=\frac{q}{2\pi}\ln(z-z_0)
\tag{5.42b}
\end{equation*}
点源流动的流线谱和等势线如图5.5所示.等势线是圆周线，流函数等值线是平面径向直线.

\textbf{(3)平面偶极子}

两无限长直线点源相距$\delta l$，线源强度分别为$q$（位于$z=-\delta l$）和$-q$（位于$z=0$），当$\delta l\to 0$时有$\displaystyle\lim_{\delta l\to 0}q\delta l=m$，则称这一对直线点源为平面偶极子.平面偶极子的复势可导出如下：
\begin{equation*}
W(z)=\lim_{\delta l\to 0}\frac{q}{2\pi}\big[\ln(z+\delta l)-\ln z\big]=\lim_{\delta l\to 0}\frac{q\delta l}{2\pi}\,\frac{\mathrm{d}}{\mathrm{d}z}\ln(z)=\lim_{\delta l\to 0}\frac{q\delta l}{2\pi z}=\frac{m}{2\pi z}
\tag{5.43}
\end{equation*}
平面偶极子的势函数和流函数可由$W(z)$的实部和虚部求出：
\begin{equation*}
\Phi=\frac{m}{2\pi}\,\frac{x}{x^{2}+y^{2}}
\tag{5.44}
\end{equation*}
\begin{equation*}
\Psi=\frac{-m}{2\pi}\,\frac{y}{x^{2}+y^{2}}
\tag{5.45}
\end{equation*}
平面偶极子的流线和等势线如图5.6所示.

\textbf{(4)点涡}

无限长孤立直涡线周围的理想不可压缩流场也是无旋的，并且是垂直于涡线的平面流动，这种平面流动称为点涡.它的复势用以下方法求出.

设直涡线通过坐标原点，并和$z$轴重合，在垂直于线涡的平面内作任意半径的圆，由于流动的对称性，速度场$v_r=0$，$v_\theta=v_\theta(r)$，沿圆周线作线积分，其环量值应等于线涡强度$\Gamma$，即$2\pi rv_\theta=\Gamma$，或$v_\theta=\Gamma/2\pi r$.由速度势定义：$\Phi=\displaystyle\int\boldsymbol{U}\cdot\mathrm{d}\boldsymbol{x}=\int v_\theta r\,\mathrm{d}\theta$，得
\begin{equation*}
\Phi=\frac{\Gamma}{2\pi}\,\theta=\frac{\Gamma}{2\pi}\arctan\frac{y}{x}
\tag{5.46a}
\end{equation*}
由速度势和流函数关系式(5.31)，可求出流函数为
\begin{equation*}
\Psi=-\frac{\Gamma}{2\pi}\ln r
\tag{5.46b}
\end{equation*}
将势函数和流函数代入复势公式，得平面不可压缩点涡的复势
\begin{equation*}
W=\Phi+\mathrm{i}\Psi=\frac{\Gamma}{2\pi\mathrm{i}}\ln z
\tag{5.47}
\end{equation*}
点涡的平面流场流线和等势线示于图5.7.

\textbf{(5)角域流动}

最后考察复势为幂函数的流型：
\begin{equation*}
W=Az^{n}
\tag{5.48}
\end{equation*}
其中$A$是实数，$n$是正实数.用极坐标表示复数$z=r\exp(\mathrm{i}\theta)$，则该复势可写成
\[
W=Ar^{n}\exp(\mathrm{i}n\theta)=A\,(r^{n}\cos n\theta+\mathrm{i}\,r^{n}\sin n\theta)
\]
将实部和虚部分离，得势函数和流函数分别为
\begin{equation*}
\Phi=Ar^{n}\cos n\theta
\tag{5.49}
\end{equation*}
\begin{equation*}
\Psi=Ar^{n}\sin n\theta
\tag{5.50}
\end{equation*}
绕角流动的速度场为
\begin{equation*}
\bar{U}=\frac{\mathrm{d}W}{\mathrm{d}z}=Anz^{n-1}
\tag{5.51}
\end{equation*}
下面我们考察不同$n$值所对应的流动形态.

\textcircled{1} $n=1$，则$W=Az$，这是理想流体在平面上的均匀流动（图5.8(a)）.

\textcircled{2} $n=1/2$，即$W=A\sqrt{z}=A(\sqrt{r}\cos(\theta/2)+\mathrm{i}\sqrt{r}\sin(\theta/2))$，当$\theta=0$和$\theta=2\pi$时$\Psi=0$，它们是两条流线.它的流线谱如图5.8(b)所示，因此这种流动是绕$2\pi$角流动.

\textcircled{3} $1>n>1/2$，由$W=Az^{n}=A(r^{n}\cos n\theta+\mathrm{i}\,r^{n}\sin n\theta)$，当$\theta=0$和$\theta=\pi/n>\pi$时，$\Psi=0$，它们是两条流线.其流线谱如图5.8(c)所示，这是绕外角（$>\pi$）的流动.

\textcircled{4} $n>1$，这时，当$\theta=0$和$\theta=\pi/n<\pi$时，$\Psi=0$，它们是两条流线，其流线谱如图5.8(d)所示，这是一种绕内角（$<\pi$）流动.

以上四种情况基本上可以分为三类：$n=1$的平行流动；$n>1$对应于绕内角流动（绕角小于$\pi$）；$n<1$对应于绕外角流动（绕角大于$\pi$）.后两类流场有以下不同特点：

（i）$n>1$，绕内角流动
\[
\bar{U}(z=0)=\big(Anz^{n-1}\big)_{z=0}=0\,,\qquad \bar{U}(z\to\infty)=\big(Anz^{n-1}\big)_{z\to\infty}\to\infty
\]
即绕内角流动时，角点是驻点（速度等于零）；而在无穷远处速度为无穷大.

（ii）$n<1$，绕外角流动
\[
\bar{U}(z=0)=\big(Anz^{n-1}\big)_{z=0}\to\infty\,,\qquad \bar{U}(z\to\infty)=\big(Anz^{n-1}\big)_{z\to\infty}=0
\]
绕外角的无旋流动中，无穷远处流速为零，但角点的速度是无穷大.

实际上不存在完全的绕角流动.但是如果物面上局部有角点（一般是绕外角），则角点附近速度势可近似为绕角流动，如果将绕流的复势在角点处作级数展开，则流动复势在角点附近的渐近展开式应为：$\displaystyle\lim_{z\to z_0}W(z)\approx A(z-z_0)^{n}$.于是，由绕角域无旋流动的性质可知，无旋流动绕外角点附近速度很大（理论上是无穷大）.

作为复势方法的应用，下面先用复势的基本解通过叠加方法来构造绕圆柱的无旋流动.

\section*{5.4  不可压缩流体绕圆柱的定常无旋流动}\addcontentsline{toc}{section}{5.4 不可压缩流体绕圆柱的定常无旋流动}

\subsection*{1. 绕圆柱无环量流动的复势}

和圆球的不可压缩无旋绕流解类似，选取均匀流（$U_\infty$为实数）和平面偶极子复势的叠加作为一种绕封闭物面的平面无旋流动，这种流动的复势表达式为
\begin{equation*}
W(z)=U_\infty z+\frac{m}{2\pi z}
\tag{5.52a}
\end{equation*}
它也可写作
\begin{equation*}
W(z)=U_\infty x+\frac{mx}{2\pi(x^{2}+y^{2})}+\mathrm{i}\Big(U_\infty y-\frac{my}{2\pi(x^{2}+y^{2})}\Big)
\tag{5.52b}
\end{equation*}
该复势满足以下条件：

（1）无穷远处为均匀来流.
\[
\Big(\frac{\mathrm{d}W}{\mathrm{d}z}\Big)_{z\to\infty}=\Big(U_\infty-\frac{m}{2\pi z^{2}}\Big)_{z\to\infty}=U_\infty
\]
（2）绕圆周线的环量等于零.应用解析函数的留数定理可证明该流动绕圆周线的环量等于零，即
\[
\Gamma=\oint\mathrm{d}W=\oint\frac{\mathrm{d}W}{\mathrm{d}z}\,\mathrm{d}z=\oint_L\Big(U_\infty-\frac{m}{2\pi z^{2}}\Big)\mathrm{d}z=0
\]
（3）圆周线是流线.取复势的虚部为零，得
\[
U_\infty y-\frac{my}{2\pi(x^{2}+y^{2})}=0
\]
该代数方程的解为
\[
y=0\qquad\text{和}\qquad x^{2}+y^{2}=\frac{m}{2\pi U_\infty}
\]
上式表示流函数等于零的流线由$x$轴（$y=0$）和圆周线组成，因而该复势满足圆周线上的不可穿透条件.圆周半径：$a=\sqrt{\dfrac{m}{2\pi U_\infty}}$，或$m=2\pi U_\infty a^{2}$，将它代入式(5.52b)，可得均匀来流绕半径为$a$的圆柱无环量流动的复势：
\begin{equation*}
W(z)=U_\infty z+\frac{a^{2}U_\infty}{z}
\tag{5.53}
\end{equation*}

\subsection*{2. 绕圆柱体无环量流动的速度场和压强分布}

由复势可以求出绕圆柱体不可压缩流体无环量流动的势函数和流函数分别为
\begin{equation*}
\Phi=U_\infty x+\frac{a^{2}U_\infty x}{x^{2}+y^{2}}=U_\infty r\cos\theta+\frac{a^{2}U_\infty\cos\theta}{r}
\tag{5.54}
\end{equation*}
\begin{equation*}
\Psi=U_\infty y-\frac{a^{2}U_\infty y}{x^{2}+y^{2}}=U_\infty r\sin\theta-\frac{a^{2}U_\infty\sin\theta}{r}
\tag{5.55}
\end{equation*}
流线谱为$\Psi=\mathrm{const.}$：
\begin{equation*}
U_\infty y-\frac{a^{2}U_\infty y}{x^{2}+y^{2}}=\mathrm{const.}
\tag{5.56}
\end{equation*}
这是一族三次曲线，并示于图5.9.

绕圆柱体不可压缩无旋流动的速度场为
\begin{equation*}
u=\frac{\partial\Phi}{\partial x}=U_\infty+a^{2}U_\infty\,\frac{y^{2}-x^{2}}{(x^{2}+y^{2})^{2}}
\tag{5.57a}
\end{equation*}
\begin{equation*}
v=\frac{\partial\Phi}{\partial y}=-a^{2}U_\infty\,\frac{2xy}{(x^{2}+y^{2})^{2}}
\tag{5.57b}
\end{equation*}
或用极坐标表示为
\begin{equation*}
v_r=U_\infty\Big(1-\frac{a^{2}}{r^{2}}\Big)\cos\theta
\tag{5.57c}
\end{equation*}
\begin{equation*}
v_\theta=-U_\infty\Big(1+\frac{a^{2}}{r^{2}}\Big)\sin\theta
\tag{5.57d}
\end{equation*}
由解出的速度场可以得到该流场的主要特性如下.

\textbf{(1)驻点}

令复速度等于零，由$\mathrm{d}W/\mathrm{d}z=U_\infty-a^{2}U_\infty/z^{2}=0$，解出$z=\pm a$，即在物面上$z=a$和$z=-a$的两点上流体速度等于零.$z=-a$在圆周线的迎风面，称前驻点；$z=a$在圆周线的背风面，称后驻点.

\textbf{(2)速度绝对值最大点}

第4章曾经证明，不可压缩无旋流动的速度绝对值最大点在物面上.将$r=a$代入速度公式(5.57c)和(5.57d)，得物面速度：
\[
v_r=0\,,\qquad v_\theta=-2U_\infty\sin\theta
\]
由上式可得速度绝对值最大点位于$r=a$，$\theta=\pm\pi/2$，或$z=\pm\mathrm{i}a$.最大速度值
\begin{equation*}
U_{\max}=2U_\infty
\tag{5.58}
\end{equation*}
\textbf{(3)流场压力分布和物面上压强的合力}

将速度代入伯努利公式，得
\[
p=p_\infty+\frac{\rho U_\infty^{2}}{2}-\frac{\rho}{2}\,(u^{2}+v^{2})=p_\infty+\frac{\rho U_\infty^{2}}{2}-\frac{\rho}{2}\,(v_r^{2}+v_\theta^{2})
\]
由此可得物面上（$r=a$）压强分布
\begin{equation*}
p=p_\infty+\frac{\rho U_\infty^{2}}{2}-\frac{\rho}{2}\,\big(4U_\infty^{2}\sin^{2}\theta\big)=p_\infty+\frac{\rho U_\infty^{2}}{2}\,\big(1-4\sin^{2}\theta\big)
\tag{5.59}
\end{equation*}
物面上压强常用无量纲压强系数表示，压强系数的定义为
\begin{equation*}
C_p=\frac{p-p_\infty}{\rho U_\infty^{2}\big/2}
\tag{5.60}
\end{equation*}
因此不可压缩无旋绕流在圆柱面上的压强系数
\begin{equation*}
C_p=1-4\sin^{2}\theta
\tag{5.61}
\end{equation*}
压强最大值在$\theta=0,\pi$，即前后驻点处，驻点压强系数$C_{p0}=1$；压强最小值在$\theta=\pm\pi/2$处，最小压强系数$C_{p\min}=-3$.

将物面压强代入合力公式，可得圆周线上的合力（圆柱轴向单位长度上的作用力）：
\[
\boldsymbol{F}=-\int_0^{2\pi}p\boldsymbol{n}\,\mathrm{d}S=-\int_0^{2\pi}p\boldsymbol{e}_r a\,\mathrm{d}\theta=-\int_0^{2\pi}a\bigg[p_\infty+\frac{\rho U_\infty^{2}}{2}\,\big(1-4\sin^{2}\theta\big)\bigg](\cos\theta\,\boldsymbol{e}_x+\sin\theta\,\boldsymbol{e}_y)\,\mathrm{d}\theta
\]
不难验证，上述积分结果等于零.就是说：不可压缩流体绕圆柱的无环量无旋流动作用在圆柱上的合力为零.图5.10给出圆柱面上压强系数的分布.

在图5.10中，同时给出圆周线上计算得到的压力系数及实际压力系数的分布，图中右边还给出不同流速下的流线图形.与圆球绕流类似，实际流动中圆柱体后部压强较理论值低，它不能回升到理论最大值，这是由于流体粘性和流动分离的缘故.图中有三条典型的压强分布曲线和典型的绕流流线图.曲线Ⅰ对应于理想无环量绕流；曲线Ⅱ对应于有分离绕流，圆周上分离点的角度小于$90^{\circ}$（以前驻点角度为零计算）；曲线Ⅲ对应于分离点角度大于$90^{\circ}$的有分离绕流.关于有分离流动的情况将在粘性流动中讨论.从图5.10可得出以下结论：①在分离以前理想绕流的压强分布和实际情况符合较好；②由于流动分离，圆柱体后驻点附近压强低于迎风面的压强，因而圆柱体上受到流动阻力，但升力仍为零.

\subsection*{3. 绕圆柱体有环量的无旋流场}

\textbf{(1)绕圆柱体有环量的无旋流动的复势}

复连通域中的无旋流动可能存在绕物面的环量$\Gamma$，本节研究绕圆柱体有环量的无旋流动.前面已经给出过点涡解，它的流线是一族绕点涡的同心圆，因此在圆柱的中心设置点涡仍满足圆柱面上的不可穿透边界条件.在绕圆柱的无环量解中叠加一点涡复势，得如下不可压缩无旋流动复势：
\begin{equation*}
W(z)=U_\infty z+\frac{U_\infty a^{2}}{z}+\frac{\Gamma}{2\pi\mathrm{i}}\ln z
\tag{5.62}
\end{equation*}
可验证该复势满足给定的边界条件.

\textbf{无穷远处来流条件}
\[
\Big(\frac{\mathrm{d}W}{\mathrm{d}z}\Big)_\infty=\Big(U_\infty-\frac{U_\infty a^{2}}{z^{2}}+\frac{\Gamma}{2\pi\mathrm{i}z}\Big)_{z\to\infty}=U_\infty
\]
\textbf{环量条件}

应用解析函数的留数定理可证明，式(5.62)的复势绕圆周线的环量：
\[
\oint\mathrm{d}W=\oint\Big(U_\infty-\frac{U_\infty a^{2}}{z^{2}}+\frac{\Gamma}{2\pi\mathrm{i}z}\Big)\mathrm{d}z=\Gamma
\]
\textbf{物面不可穿透条件}

圆周线上流函数值
\[
\Psi\big|_{z=a\exp(\mathrm{i}\theta)}=\big[\mathrm{Im}W(z)\big]_{z=a\exp(\mathrm{i}\theta)}=-\frac{\Gamma}{2\pi}\ln a=\mathrm{const.}
\]
因此，该复势满足物面不可穿透条件.于是，式(5.62)表示的复势是均匀来流绕圆柱体有环量的不可压缩流体无旋流动解.

\textbf{(2)绕圆柱体有环量无旋流动的速度场}

由复势可计算复速度
\begin{equation*}
\frac{\mathrm{d}W}{\mathrm{d}z}=U_\infty-\frac{U_\infty a^{2}}{z^{2}}+\frac{\Gamma}{2\pi\mathrm{i}z}
\tag{5.63}
\end{equation*}
由此可导出速度场为
\[
u=U_\infty+\frac{U_\infty a^{2}\,(y^{2}-x^{2})}{(x^{2}+y^{2})^{2}}-\frac{\Gamma y}{2\pi(x^{2}+y^{2})}\,,\qquad v=-\frac{2U_\infty a^{2}xy}{(x^{2}+y^{2})^{2}}+\frac{\Gamma x}{2\pi(x^{2}+y^{2})}
\]
或用极坐标表示为
\[
v_r=U_\infty\Big(1-\frac{a^{2}}{r^{2}}\Big)\cos\theta\,,\qquad v_\theta=-U_\infty\Big(1+\frac{a^{2}}{r^{2}}\Big)\sin\theta+\frac{\Gamma}{2\pi r}
\]
通过复势和复速度可以分析绕圆柱体的不可压缩有环量无旋流动的主要性质如下.

\textcircled{1} 驻点位置

令$\mathrm{d}W/\mathrm{d}z=0$，解出
\begin{equation*}
z=\frac{\mathrm{i}\Gamma}{4\pi U_\infty}\pm\sqrt{a^{2}-\frac{\Gamma^{2}}{16\pi^{2}U_\infty^{2}}}
\tag{5.64}
\end{equation*}
我们来讨论以下三种情况.

（i）$a^{2}-\dfrac{\Gamma^{2}}{16\pi^{2}U_\infty^{2}}>0$或$\dfrac{|\Gamma|}{4\pi U_\infty a}<1$.这时式(5.64)有两个解：
\[
z_{1,2}=\frac{\mathrm{i}\Gamma}{4\pi U_\infty}\pm\sqrt{a^{2}-\frac{\Gamma^{2}}{16\pi^{2}U_\infty^{2}}}
\]
两个驻点的坐标分别为
\[
\Big(\sqrt{a^{2}-\frac{\Gamma^{2}}{16\pi^{2}U_\infty^{2}}}\,,\ \frac{\Gamma}{4\pi U_\infty}\Big)\qquad\text{和}\qquad\Big(-\sqrt{a^{2}-\frac{\Gamma^{2}}{16\pi^{2}U_\infty^{2}}}\,,\ \frac{\Gamma}{4\pi U_\infty}\Big)
\]
两个驻点$z_1,z_2$的模$|z_1|=|z_2|=a$，就是说，这两驻点都在圆周线上（见图5.11(a)），并相对于$y$轴对称.

（ii）$\Gamma=\pm 4\pi U_\infty a$

这时只有一个驻点，当$\Gamma=4\pi U_\infty a$时，$z=\Gamma\mathrm{i}/4\pi U_\infty=a\mathrm{i}$，即正环量时，驻点在正$y$轴上；当$\Gamma=-4\pi U_\infty a$时，$z=\mathrm{i}\Gamma/4\pi U_\infty=-a\mathrm{i}$，即负环量时，驻点在圆柱的下方（图5.11(b)）.

（iii）$\Gamma^{2}>16\pi^{2}a^{2}U_\infty^{2}$或$a^{2}-\dfrac{\Gamma^{2}}{16\pi^{2}U_\infty^{2}}<0$，这时驻点位于：
\[
z=\Big(\frac{\Gamma}{4\pi U_\infty}\pm\sqrt{\frac{\Gamma^{2}}{16\pi^{2}U_\infty^{2}}-a^{2}}\Big)\mathrm{i}
\]
这时形式上仍有两个驻点.其中一个在圆外$|z|>a$，而另一个在圆内$|z|<a$（流场外），所以绕圆周流场中只有一个圆外的驻点，
\[
z_1=\Big(\frac{\Gamma}{4\pi U_\infty}\pm\sqrt{\frac{\Gamma^{2}}{16\pi^{2}U_\infty^{2}}-a^{2}}\Big)\mathrm{i}
\]
式中“$+$”对应于$\Gamma>0$，驻点位于圆周线上面；“$-$”对应于$\Gamma<0$，驻点位于圆周线下面.以上三种情况示于图5.11.

\textcircled{2} 圆周线（或圆柱面）上的速度分布

将圆柱周线方程$r=a$代入复速度公式，得
\[
v_r=0\,,\qquad v_\theta=-2U_\infty\sin\theta+\Gamma\big/(2\pi a)
\]
因此
\[
U^{2}=\Big(2U_\infty\sin\theta-\frac{\Gamma}{2\pi a}\Big)^{2}
\]
因$\sin\theta=\sin(\pi-\theta)$，故速度分布对于$y$轴对称；而$\sin(-\theta)=-\sin\theta$，因此不可压缩流体绕圆柱的有环量流动的速度场对$x$轴反对称.

\textcircled{3} 圆柱面上的压强分布

将圆周线速度分布代入伯努利公式，可得圆柱面上压强分布：
\[
p=p_\infty+\frac{\rho U_\infty^{2}}{2}-\frac{\rho U^{2}}{2}=p_\infty+\frac{\rho U_\infty^{2}}{2}\bigg[1-\Big(4\sin^{2}\theta-\frac{2\Gamma\sin\theta}{\pi aU_\infty}+\frac{\Gamma^{2}}{4\pi^{2}a^{2}U_\infty^{2}}\Big)\bigg]
\]
由于$\sin\theta=\sin(\pi-\theta)$和$\sin(-\theta)=-\sin\theta$，压强分布对$y$轴对称，而对$x$轴反对称.这意味着，这种有环量绕流作用在圆柱上的阻力（来流方向的合力）仍然等于零，而垂直于来流方向的合力将不等于零.

\textcircled{4} 圆柱周线上的合力

将压强分布代入合力公式，得
\begin{align*}
\boldsymbol{F}&=-\int_0^{2\pi}p\boldsymbol{n}a\,\mathrm{d}\theta=-a\int_0^{2\pi}p\,(\cos\theta\,\boldsymbol{e}_x+\sin\theta\,\boldsymbol{e}_y)\,\mathrm{d}\theta\\
&=-\int_0^{2\pi}a\,\bigg\{p_\infty+\frac{\rho U_\infty^{2}}{2}\bigg[1-\Big(4\sin^{2}\theta-\frac{2\Gamma\sin\theta}{\pi aU_\infty}+\frac{\Gamma^{2}}{4\pi^{2}a^{2}U_\infty^{2}}\Big)\bigg]\bigg\}\cdot(\cos\theta\boldsymbol{e}_x+\sin\theta\boldsymbol{e}_y)\,\mathrm{d}\theta
\end{align*}
积分后得
\begin{equation*}
\boldsymbol{F}=-\rho U_\infty\Gamma\,\boldsymbol{e}_y
\tag{5.65}
\end{equation*}
这说明绕圆柱的有环量无旋流动中，在圆柱上有垂直于流动方向的作用力，并称之为升力.当来流$U_\infty$为正$x$轴方向时，正环量（逆时针）产生的升力指向负$y$轴方向；而负环量（顺时针）则产生正$y$方向的升力.如环量用涡线方向表示：$\boldsymbol{\Gamma}=\Gamma\boldsymbol{k}$；来流也用向量表示：$\boldsymbol{U}_\infty=U_\infty\boldsymbol{i}$，则升力公式可表示为
\begin{equation*}
\boldsymbol{F}=\rho\,\boldsymbol{U}_\infty\times\boldsymbol{\Gamma}
\tag{5.66}
\end{equation*}
在不可压缩流体绕圆柱有环量无旋流动情况下，圆柱周线上升力产生的原因可由无环量绕流和点涡叠加的流场来给予解释.以图5.11(a)为例（负环量，即点涡流场是顺时针方向），由于圆柱面下方由环量产生的速度和绕圆柱无环量绕流速度是反向的，因而速度值减小，根据伯努利公式压强值较无环流时增大；反之在圆柱面上方，速度值增大，压强减小，因此负环量（顺时针）绕圆柱流动的情况必产生一向上的合力.正环量（逆时针）的情况下产生向下的升力可作同样的解释.

圆柱体有环量绕流产生升力的现象称为马格努斯（Magnus）效应.马格努斯曾设想利用旋转圆柱产生环量，从而在风速下产生横向力，以取代风帆.但是实际情况下马格努斯得到的升力远小于理想绕流的$\rho U_\infty\Gamma$，这是由于圆柱体外的真实绕流不能保持理想的无旋流动状态，例如柱体的后部必有流动分离现象，如图5.10所示.

\section*{5.5  解平面不可压缩无旋绕流的保角映像法}\addcontentsline{toc}{section}{5.5 解平面不可压缩无旋绕流的保角映像法}

下面介绍用复变函数的保角映像法求解不可压缩平面无旋绕流问题.

\subsection*{1. 保角映像法的基本原理}

保角映像法的基本思想是利用解析变换来寻求满足边界条件的复势，具体做法如下.

\textbf{(1)解析变换的保角映射}

数学上，解析函数$z=Z(\zeta)$是一种几何变换，它可以将$\zeta$平面上的曲线变换到$z$平面上另一几何曲线（或反之），所以解析变换是建立两个平面之间几何关系的一种方法，举例如下.

设有$z=x+\mathrm{i}y$平面上的几何曲线$K$：$z(t)=x(t)+\mathrm{i}y(t)$，利用解析函数$\zeta=\zeta(z)$变换到$\zeta=\xi+\mathrm{i}\eta$平面上的曲线$K^{*}$：$\zeta=\zeta(t)=\xi(t)+\mathrm{i}\eta(t)$.具体做法是将曲线方程$z(t)=x(t)+\mathrm{i}y(t)$代入解析函数式，得$\zeta(t)=\zeta(z(t))=\xi(t)+\mathrm{i}\eta(t)$.我们称$z$平面为原平面或物理平面；$\zeta=\xi+\mathrm{i}\eta$为变换平面，简称$\zeta$平面.由于解析函数的性质，这种变换是一一对应和可导的，因此$z$平面上的封闭曲线变换到$\zeta$平面也是封闭曲线，并且变换式$\zeta=\zeta(z)$必有逆变换$z=z(\zeta)$.解析变换有一个重要的性质：任意两曲线的夹角在解析变换下保持不变，如：两正交的曲线在解析函数变换过程中仍然保持正交性.这是保角变换名词的来源.

\textbf{(2)由解析变换构造新的复势}

利用解析变换的性质，可以构造不可压缩平面无旋流动的不同复势.如果在$z$平面上有一复势$W(z)$，它表示某种平面不可压缩无旋流动.现有一解析变换$z=Z(\zeta)$，它在$z$平面和$\zeta$平面间建立一一对应关系.将$z=Z(\zeta)$代入$W(z)$，构成另一复函数$W^{*}(\zeta)=W(Z(\zeta))$.因为解析函数的复合函数仍然是解析函数，$W^{*}(\zeta)$表示$\zeta$平面上的某种复势.反之，已知$\zeta$平面的复势$W^{*}(\zeta)$，通过解析变换$\zeta=\zeta(z)$可以获得$z$平面上的复势$W(z)$.下面具体说明这种变换的流动性质.

已知绕圆周线的不可压缩平面无旋流动的复势，设该流动的自变量为$\zeta$，绕流的复势（见式(5.62)）：
\[
W^{*}(\zeta)=U_\infty^{*}\zeta+\frac{U_\infty^{*}a^{2}}{\zeta}+\frac{\Gamma}{2\pi\mathrm{i}}\ln\zeta
\]
通过$W^{*}(\zeta)$，我们可以构造新的无旋绕流复势.具体做法是：给定一个在圆外解析函数$\zeta=\zeta(z)$（它的反函数$z=Z(\zeta)$，也是解析函数），并要求它具有以下性质.

\textcircled{1} 它将$\zeta$平面上的圆周线$K^{*}\,(\zeta=a\exp(\mathrm{i}\theta))$变换成$z$平面上的某一封闭曲线$K\,(z=z(\theta))$.

\textcircled{2} 无穷远点对应：
\begin{equation*}
z(\infty)=\infty
\tag{5.67}
\end{equation*}
\textcircled{3} 无穷远处变换函数的导数是有限复常数
\begin{equation*}
\frac{\mathrm{d}z}{\mathrm{d}\zeta}(\infty)=m\exp(\mathrm{i}\alpha)<\infty\,,\qquad \frac{\mathrm{d}\zeta}{\mathrm{d}z}(\infty)=\frac{1}{m}\exp(-\mathrm{i}\alpha)<\infty
\tag{5.68}
\end{equation*}
式中$m$是实常数.在复变函数理论中已经证明这种解析变换存在且唯一.

将$\zeta=\zeta(z)$代入绕圆周线的复势，将对应一个新的复函数如下：
\begin{equation*}
W(z)=W^{*}(\zeta(z))=U_\infty^{*}\zeta(z)+\frac{U_\infty^{*}a^{2}}{\zeta(z)}+\frac{\Gamma}{2\pi\mathrm{i}}\ln\zeta(z)
\tag{5.69}
\end{equation*}
因为解析函数的复合函数仍是解析函数，所以新的复函数也是解析函数，因而它是某一流动的复势，并具有以下性质.

\textcircled{1} 变换后的周线$K\,(z=z(\theta))$是流线，即$\mathrm{Im}[W(z(\theta))]=$实常数.

因为在$\zeta$平面上圆周线$K^{*}\,(\zeta=a\exp(\mathrm{i}\theta))$是流线，即$\mathrm{Im}[W^{*}(\zeta)]_{\zeta=a\exp(\mathrm{i}\theta)}=$实常数，而周线$K\,(z=z(\theta))$和圆周线$K^{*}\,(\zeta=a\exp(\mathrm{i}\theta))$是一一对应的，在对应点上的复势值相等（复势变换式(5.69)），所以，在封闭周线$K\,(z(\theta))$上复势的虚部
\[
\mathrm{Im}\{W[z(\theta)]\}=\mathrm{Im}\big[W^{*}(\zeta)\big]_{\zeta=a\exp(\mathrm{i}\theta)}=\text{实常数}
\]
因此周线$K\,(z=z(\theta))$也是流线.

\textcircled{2} 两流场的无穷远点相对应（式(5.67)），也就是说无界的绕圆周线的流动问题对应一个新的无界绕流问题.

\textcircled{3} 新复势无穷远的来流速度大小放大$1/m$倍，而速度方向角增加$\alpha$.

因为通过复合函数求导数可求得新复势的复速度为
\begin{equation*}
\frac{\mathrm{d}W(z)}{\mathrm{d}z}=\frac{\mathrm{d}W^{*}(\zeta)}{\mathrm{d}\zeta}\,\frac{\mathrm{d}\zeta}{\mathrm{d}z}
\tag{5.70}
\end{equation*}
对应$z$平面上无穷远处来流的复速度为
\begin{align*}
U_\infty+\mathrm{i}V_\infty&=|U|_\infty\exp(\mathrm{i}\beta)=\bigg[\frac{\mathrm{d}W(z)}{\mathrm{d}z}\bigg]_{z\to\infty}\\
&=\bigg[\frac{\mathrm{d}W^{*}(\zeta)}{\mathrm{d}\zeta}\bigg]_{\zeta\to\infty}\bigg[\frac{\mathrm{d}\zeta}{\mathrm{d}z}\bigg]_{z\to\infty}=\frac{1}{m}\,U_\infty^{*}\exp(\mathrm{i}\alpha)
\tag{5.71}
\end{align*}
上式清楚地说明$\Big|\dfrac{\mathrm{d}W}{\mathrm{d}z}\Big|_\infty=\dfrac{1}{m}U_\infty$和$\beta=\alpha$，即来流速度放大$1/m$倍，来流方向转过$\alpha$.如果要求变换式的$m=1$和$\alpha=0$，则变换前后无穷远处的速度大小相等、方向相同.

\textcircled{4} 新复势绕物面的环量和原绕圆周线的环量相等.

前面已经证明绕物面的环量等于$\Gamma=\displaystyle\oint_K\mathrm{d}W$，在对应的物面上有：$\mathrm{d}W^{*}(\zeta)=\mathrm{d}W(z)$，因此
\begin{equation*}
\Gamma=\oint_K\mathrm{d}W=\oint_{K^{*}}\mathrm{d}W^{*}
\tag{5.72}
\end{equation*}
通过以上分析可以看到，通过圆周线外的解析变换$\zeta=\zeta(z)$，可以将不可压缩流体绕圆柱的无旋流动变换成来流速度大小和环量都不变的绕封闭型线$z=z(\theta)$的无旋流动.上述保角映射或解析变换构造复势的原理示于图5.12.

根据上面给出的解析变换方法，可以用绕圆柱无旋流动的复势作为基本解，通过解析变换，得到各种平面无旋流动的复势.下面给出正反两类问题的解法.

\textcircled{1} \textbf{反问题方法}：以变换平面上绕圆柱有环量无旋流动的复势为基本解，给定变换函数$\zeta(z)$，求物理平面对应的流动.已知不可压缩流体绕圆柱无旋流动的复势为
\[
W^{*}(\zeta)=U_\infty^{*}\zeta+\frac{U_\infty^{*}a^{2}}{\zeta}+\frac{\Gamma}{2\pi\mathrm{i}}\ln\zeta
\]
给定满足式(5.67)、式(5.68)的变换函数$\zeta=\zeta(z)$（反函数为$z=z(\zeta)$），将它代入基本解，得物理平面绕新型线$z=z(\theta)$的无旋流动的复势$W(z)$如下：
\begin{equation*}
W(z)=W^{*}(\zeta(z))=U_\infty^{*}\zeta(z)+\frac{U_\infty^{*}a^{2}}{\zeta(z)}+\frac{\Gamma}{2\pi\mathrm{i}}\ln\zeta(z)
\tag{5.73}
\end{equation*}
新型线的几何方程为$z=z(a\exp(\mathrm{i}\theta))$.

\textcircled{2} \textbf{正问题方法}：给定物理平面的物面型线方程：$z=z(t)$，求出满足式(5.67)、式(5.68)的解析变换$\zeta=\zeta(z)$（或$z=z(\zeta)$），将给定的物面型线解析变换成半径为$a$的圆，即$\zeta(z(t))=a\exp(\mathrm{i}t)$，然后将变换式$\zeta=\zeta(z)$代入圆柱绕流的基本解，就可得到给定物形的无旋绕流的复势.

反问题方法是先给出解析变换函数，然后确定对应这种变换的绕流物面型线.正问题方法是先给定绕流的物面型线，然后确定满足式(5.67)、式(5.68)的解析变换式.

下面我们用实例具体说明保角映射解法.

\subsection*{2. 简单的解析函数及反问题}

\textbf{(1)线性函数}
\begin{equation*}
\zeta=az+b=c\,(z-z_0)
\tag{5.74}
\end{equation*}
线性变换在几何上是坐标平移、放大和旋转.$z_0$是坐标原点的平移.$c=m\exp(-\mathrm{i}\alpha)$是坐标的放大（放大倍数$m$）和顺时针旋转（角度$\alpha$），不难验证，它具有保角性.

将变换式代入绕圆柱无旋流动的复势$W^{*}(\zeta)=U_\infty^{*}\zeta+\dfrac{U_\infty^{*}a^{2}}{\zeta}+\dfrac{\Gamma}{2\pi\mathrm{i}}\ln\zeta$，得到新复势为
\begin{equation*}
W(z)=cU_\infty^{*}(z-z_0)+\frac{U_\infty^{*}a^{2}}{c(z-z_0)}+\frac{\Gamma}{2\pi\mathrm{i}}\ln\big(c(z-z_0)\big)
\tag{5.75a}
\end{equation*}
下面来考察新复势所表示的绕流及其性质.将圆周线方程$\zeta=a\exp(\mathrm{i}\theta)$代入$\zeta=c(z-z_0)$，新复势的绕流型线$K$为
\[
z-z_0=\frac{a}{c}\exp(\mathrm{i}\theta)=\frac{a}{m}\exp(\mathrm{i}(\theta+\alpha))
\]
由上式很容易证明：$|z(\theta)-z_0|=\dfrac{a}{m}$.这表明$\zeta$平面上的圆周线在$z$平面仍然是一个圆，它的圆心在$z_0$、半径为$a/m$.

新的复势在无穷远处的流速为
\[
U_\infty-\mathrm{i}V_\infty=\Big(\frac{\mathrm{d}W}{\mathrm{d}z}\Big)_\infty=\Big(\frac{\mathrm{d}W^{*}}{\mathrm{d}\zeta}\Big)\frac{\mathrm{d}\zeta}{\mathrm{d}z}=mU_\infty^{*}\exp(-\mathrm{i}\alpha)
\]
即对应无穷远来流速度大小为$mU_\infty^{*}$，来流与水平轴交角为$\alpha$.线性变换生成的绕流图形示于图5.13.我们可以清楚地看到，线性变换$\zeta=az+b=c(z-z_0)$将绕圆周线的无旋流动变换成另一个绕圆周线的无旋流动，新的圆柱绕流的圆心在$z_0$，圆的半径等于$a/m$，来流速度等于$mU_\infty^{*}$，来流的方向和$x$轴夹角为$\alpha$.设$m=1$，并将$\alpha,z_0$等代入式(5.75)，我们就得到绕半径为$a$、圆心位于$z_0$、来流速度等于$U_\infty$、来流速度和$x$轴夹角为$\alpha$的不可压缩无旋流动的复势：
\begin{equation*}
W(z)=U_\infty\exp(-\mathrm{i}\alpha)(z-z_0)+\frac{U_\infty\exp(\mathrm{i}\alpha)a^{2}}{z-z_0}+\frac{\Gamma}{2\pi\mathrm{i}}\ln(z-z_0)-\frac{\Gamma\alpha}{2\pi}
\tag{5.75b}
\end{equation*}
式(5.75b)是不可压缩流体绕圆柱无旋流动复势的一般公式，式中最后一项是常数，它不影响速度场的计算，因此可以略去.

\textbf{(2)幂函数}
\begin{equation*}
z=A\zeta^{n}
\tag{5.76}
\end{equation*}
式中$A$和$n$都是实数.令$z=r\exp(\mathrm{i}\theta)$，$\zeta=t\exp(\mathrm{i}\phi)=A^{-1/n}r^{1/n}\,(\cos(\theta/n)+\mathrm{i}\sin(\theta/n))$.幂函数将$z$平面上的射线（$\theta=C$）变换到$\zeta$平面上的射线（$\phi=C/n$）.$n=1$是线性变换，是式(5.74)的特例，它将$z$平面上无限长直线（$\theta=0,\theta=\pi$）变换到平面上的无限长直线（$\phi=0,\phi=\pi$）.当$1/2<n<1$时，$z$平面上的无限长直线（$\theta=0$和$\theta=\pi$）变换到$\zeta$平面上的外角（$\phi=0,\phi=\pi/n>\pi$）；当$n>1$时，$z$平面上的无限长直线变换到$\zeta$平面上的内角（$\phi=0,\phi=\pi/n<\pi$），如图5.14所示.

$z$平面上的均匀流动$W(z)=Uz$，用幂函数变换到$\zeta$平面上就是前面讨论过的绕角流动：$W(\zeta)=AU\zeta^{n}$.

\textbf{(3)儒可夫斯基变换和绕椭圆柱流动}

著名儒可夫斯基变换式为
\begin{equation*}
z=\frac{1}{2}\Big(\zeta+\frac{a^{2}}{\zeta}\Big)
\tag{5.77a}
\end{equation*}
它的逆变换式为
\begin{equation*}
\zeta=z\pm\sqrt{z^{2}-a^{2}}
\tag{5.77b}
\end{equation*}
平方根函数是多值函数，取定其中一叶后，式(5.77b)中正负号就可确定，通常取正号.该变换有以下几何特性.

\textcircled{1} $\zeta=0$是奇点，它对应于$z$平面上的无穷远点.

\textcircled{2} $z(\zeta=\infty)=\infty$，即$z$平面上无穷远点对应$\zeta$平面的无穷远点.

\textcircled{3} $(\mathrm{d}z/\mathrm{d}\zeta)_\infty=1/2$，即无穷远处线段长度缩小二分之一，但线段方向不变.

\textcircled{4} $\zeta=a\exp(\mathrm{i}\theta)$的圆变换为$-a$到$a$的往返线段.将$\zeta=a\exp(\mathrm{i}\theta)$代入式(5.77a)，得
\[
z(\theta)=\frac{1}{2}\big(a\exp(\mathrm{i}\theta)+a\exp(-\mathrm{i}\theta)\big)=a\cos\theta
\]
上式表示：$\zeta$平面上$r=a$的上半圆（$0<\theta<\pi$）变换成从$z=a$到$z=-a$的线段；下半圆（$\pi<\theta<2\pi$）变换成由$z=-a$到$z=a$的线段，$\zeta$平面上的圆外域变换成$z$平面上长$2a$线段的外域.

\textcircled{5} 此变换将圆周线$\zeta=c\exp(\mathrm{i}\theta)\,(c>a)$变换成椭圆周线，将它代入式(5.77a)，得到
\[
z(\theta)=x(\theta)+\mathrm{i}y(\theta)=\frac{1}{2}\bigg[c\exp(\mathrm{i}\theta)+\frac{a^{2}}{c}\exp(-\mathrm{i}\theta)\bigg]
\]
实部和虚部分开，可得型线的参数方程：
\[
x(\theta)=\frac{1}{2}\Big(c+\frac{a^{2}}{c}\Big)\cos\theta=a_1\cos\theta\,,\qquad y(\theta)=\frac{1}{2}\Big(c-\frac{a^{2}}{c}\Big)\sin\theta=b_1\sin\theta
\]
进行简单的代数运算，可以得到$z$平面上的周线：
\[
\frac{x^{2}}{a_1^{2}}+\frac{y^{2}}{b_1^{2}}=1
\]
式中：$a_1=(c+a^{2}/c)/2$，$b_1=(c-a^{2}/c)/2$，该椭圆的焦距
\[
f=\sqrt{a_1^{2}-b_1^{2}}=a
\]
椭圆的长短轴之和等于$\zeta$平面上的圆半径：$c=a_1+b_1$.

\textcircled{6} 此变换将射线$\zeta=r\exp(\mathrm{i}\alpha)$（$\alpha$是常数，等于射线的倾斜角）变换成双曲线，将$\zeta=r\exp(\mathrm{i}\alpha)$代入式(5.77a)得
\[
z(r)=x(r)+\mathrm{i}y(r)=\frac{1}{2}\bigg[r\exp(\mathrm{i}\alpha)+\frac{a^{2}}{r}\exp(-\mathrm{i}\alpha)\bigg]
\]
将实部和虚部分开，可得型线的参数方程：
\[
x(r)=\frac{1}{2}\Big(r+\frac{a^{2}}{r}\Big)\cos\alpha\,,\qquad y(r)=\frac{1}{2}\Big(r-\frac{a^{2}}{r}\Big)\sin\alpha
\]
进行简单的代数运算消去$r$后，可以得到$z$平面上的周线：
\[
\frac{x^{2}}{\cos^{2}\alpha}-\frac{y^{2}}{\sin^{2}\alpha}=a^{2}
\]
由于在$\zeta$平面上圆周线和射线是正交的，不难验证，在$z$平面上对应的双曲线族和椭圆曲线族也是正交的（读者自行证明）.

儒可夫斯基变换的基本图形示于图5.15.

\noindent\textbf{例5.3}\quad 求来流平行于椭圆长轴的不可压缩无旋绕流的复势.

椭圆柱方程为$F(x,y)=x^{2}/a_1^{2}+y^{2}/b_1^{2}-1=0$，其来流的速度为$U_\infty\,(\alpha=0)$，已知环量为$\Gamma$，求该绕流的复势$W(z)$.

\noindent\textbf{解}\quad 利用儒可夫斯基变换
\[
\zeta=z+\sqrt{z^{2}-a^{2}}
\]
它将$\zeta$平面的圆变换为$z$平面上的椭圆，椭圆的焦距应等于变换式中的$a$，即$a=\sqrt{a_1^{2}-b_1^{2}}$，$\zeta$平面上的圆半径$R=a_1+b_1$.已知半径为$R$的圆柱绕流的复势：
\[
W^{*}(\zeta)=U_\infty^{*}\zeta+\frac{U_\infty^{*}R^{2}}{\zeta}+\frac{\Gamma}{2\pi\mathrm{i}}\ln\zeta
\]
式中
\[
U_\infty^{*}=\Big(\frac{\mathrm{d}W^{*}}{\mathrm{d}\zeta}\Big)_{\zeta\to\infty}=\Big(\frac{\mathrm{d}W}{\mathrm{d}z}\Big)_{z\to\infty}\Big(\frac{\mathrm{d}z}{\mathrm{d}\zeta}\Big)_{\zeta\to\infty}=\frac{1}{2}\Big(\frac{\mathrm{d}W}{\mathrm{d}z}\Big)_{z\to\infty}=\frac{1}{2}\,U_\infty
\]
将$U_\infty^{*}$和儒可夫斯基变换式$\zeta=z+\sqrt{z^{2}-a^{2}}$代入前式，绕椭圆的无旋流动复势为
\begin{align*}
W(z)&=\frac{U_\infty}{2}\Big(z+\sqrt{z^{2}-(a_1^{2}-b_1^{2})}\Big)+\frac{U_\infty\,(a_1+b_1)^{2}}{2\big(z+\sqrt{z^{2}-(a_1^{2}-b_1^{2})}\big)}\\
&\quad+\frac{\Gamma}{2\pi\mathrm{i}}\ln\Big(z+\sqrt{z^{2}-(a_1^{2}-b_1^{2})}\Big)
\end{align*}

\subsection*{3. 绕儒可夫斯基对称翼型流动的复势}

前面的例子说明利用解析变换可以构造新的绕流型线并给出它的复势.根据这种思想可以设计飞机或涡轮机械的翼型.典型的例子是儒可夫斯基翼型及其绕流.

\textbf{(1)儒可夫斯基对称翼型（或儒可夫斯基舵面）}

上节已经讨论过儒可夫斯基变换的一些几何变换性质.现在来进一步讨论圆心在$x$轴上的偏心圆的变换（见图5.16）.

由于$\zeta$平面上圆心在原点的圆周线经儒可夫斯基变换，在$z$平面是往返线段$AA'$，圆心位于$x$轴和与圆相切于$A''$的偏心圆$K^{*}$在儒可夫斯基变换下将变换成包含线段$AA'$的周线$K$，并且它在$A'$点和$x$轴相切.偏心圆上$F^{*}$变换到周线$K$上的前缘点$F$.于是应用儒可夫斯基变换，把偏心圆$K^{*}$变换成具有一定厚度的钝前缘
%%P4%%